%% SECTION 1: INTRODUCTION
%% ============================================================
\section{Introduction}

The conservative claim is already large: 10$\times$--30$\times$ fewer training tokens to a fixed accepted-change capability, 3$\times$--10$\times$ fewer context/reasoning/tool/retry tokens per verified change, and 3$\times$--10$\times$ lower all-in cost per verified correct change after foundry amortization. The central claim is larger: 30$\times$--150$\times$ training-token efficiency, 10$\times$--100$\times$ fewer reasoning/action/retry tokens, and 10$\times$--50$\times$ lower verified-change cost on supported canonical product work. The denominator comes first because the thesis is not ``better style.'' It is that frontier coding models are paying to learn the wrong object.

Frontier coding models are trained on a software record built by humans for humans. That record contains valuable behavior, but it is not a canonical representation of software. It is a survival archive: hundreds of languages, millions of local conventions, duplicated frameworks, inconsistent naming, weak tests, dependency sprawl, build folklore, continuous-integration (CI) dialects, repository settings, access policies, and release rituals that were never designed for autonomous agents. A coding agent pays for that archive four times: during training, during context gathering, during reasoning/tool/retry loops, and during human review.

The end-of-road story is stark. In the conservative version, canonicalization is an engineering discipline that makes coding agents cheaper and more reliable by cutting routine repository search and failed repair loops. In the central version, it becomes a new data substrate: fewer legal encodings, fewer legal edit paths, and lower tokens-to-target capability. In the aggressive version, routine product software approaches a behavior-genome regime in which most CRUD, auth, policy, migration, deployment, observability, billing, notification, workflow, and integration behavior is assembled from proof-carrying cells, while bespoke code is reserved for minimum viable novelty. The paper is ambitious because the upside is enormous; it is defensible only because every large number is tied to a named denominator and a falsification test.

The scale is enormous. GitHub reported 395 million public and open-source repositories in Octoverse 2025, roughly 230 new repositories per minute, and a major shift in the most-used language~\cite{octoverse2025}. The Stack v2 reports 67.5TB full, 32.1TB deduplicated, roughly 900B train-full tokens, over 3B files, and 658 languages~\cite{theStackV2}. StarCoder2 was trained on 3.3T--4.3T tokens across 619 programming languages~\cite{starcoder2}. Kimi K2 reports a 1T-total-parameter mixture-of-experts model with 32B active parameters and 15.5T pretraining tokens~\cite{kimiK2,kimiDocs}. Modern code-model work shows the same field-level direction: more languages, longer context, more tokens, more data curation, and more agentic training rather than a smaller human-code target~\cite{deepseekCoder,qwenCoder}.

The redundancy is also enormous. D\'{e}j\`{a}Vu analyzed 4.5 million non-fork GitHub projects and found that only 85 million of 428 million files were unique---about 70\% were clones of previously created files~\cite{dejavu2017}. Jupyter notebook code shows similar patterns: more than 70\% of code snippets were exact copies~\cite{jupyterClones2020}. Software has long been known to be highly repetitive and predictable~\cite{hindleNaturalness}.

The thesis is stronger than filtering or deduplication: the available training code should be \emph{ported} into an agent-first canonical standard before training. Filtering selects cleaner examples from the old distribution. Canonical porting changes the distribution and the action substrate. The standard constrains language roles, file grammar, naming, dependency ownership, generated boundaries, database truth, continuous integration, repository policy, build/test commands, migration policy, proof lanes, and agent-readable edit scope. Observable behavior is preserved, rejected, or tiered according to an explicit evidence envelope; arbitrary human degrees of freedom collapse only when they are not part of that envelope. The external-reader claim is precise: human repositories are improvable because they mix durable behavior with accidental representation, not because human behavior, product judgment, operational lessons, or edge cases are disposable.

This is a deliberately strong position: for every supported product-software concern there should be one governed way to express it, and for other software classes there should be a secondary governed profile rather than local preference. Humans do not merely choose different surface styles; they create mutually incompatible local dialects for routing, testing, continuous integration, authorization, secrets, migrations, feature flags, deployment, observability, and review. Those dialects may work locally, but they are training noise for agents and reasoning overhead at inference time. For supported classes, the falsifiable claim is that nearly every nontrivial human-authored repository contains reducible accidental degrees of freedom: duplicate behavior, noncanonical representation, generated-truth drift, unconstrained edit surface, missing proof receipts, or ungoverned policy.

The decisive experiment is not a better-looking repository. It is a same-model, same-lineage comparison: original human repository versus canonical repository, same issue, same behavior contract, same hidden tests, same reviewer rubric, and a complete ledger of files opened, context tokens, reasoning tokens, tool calls, invalid edits, failed proof lanes, review burden, wall time, and accepted-patch rate. If a broad model performs materially better on the canonical repository before a canonical specialist is trained, then substrate compression is real. If it does not, the thesis fails early.

This paper makes the case through nine explicit assumptions, each with falsification criteria, a compression frontier catalog, preliminary training evidence, and a paired evaluation design. The core falsifiable thesis is simple: after controlling for behavior and task, canonical proof-carrying software should lower same-model search cost now, lower tokens-to-target accepted-change capability under paired scaling curves later, and lower amortized cost per verified correct change in production field trials. If those three gates fail, no title, figure, or compression catalog rescues the claim.

%% ============================================================
